% ==============================================================================
% ĐỀ KIỂM TRA GIỮA HỌC KÌ I - MÔN TOÁN 6 - ĐỀ 2
% ==============================================================================

\noindent
\begin{minipage}[t]{6.5cm}
	\centering
	\textbf{TRƯỜNG THCS \dots\dots\dots\dots\dots\dots}\\[3pt]
	\textbf{TỔ: \dots\dots\dots\dots\dots\dots}\\[6pt]
	\textbf{(ĐỀ 2)}\\
	(\textit{Đề có 02 trang})
\end{minipage}%
\hfill
\begin{minipage}[t]{10.5cm}
	\centering
	\textbf{ĐỀ KIỂM TRA GIỮA HỌC KÌ I}\\[3pt]
	\textbf{NĂM HỌC 2026 -- 2027}\\[3pt]
	\textbf{Môn: TOÁN -- Lớp 6}\\[3pt]
	\textit{Thời gian làm bài: 90 phút (không kể thời gian giao đề)}
\end{minipage}

\vspace{10pt}

\noindent
\textbf{Họ và tên học sinh:} \dotfill{} \textbf{Lớp:} \makebox[2.5cm]{\dotfill} \textbf{Số báo danh:} \makebox[2.5cm]{\dotfill}

\vspace{10pt}

% ------------------------------------------------------------------------------
% PHẦN I: TRẮC NGHIỆM NHIỀU PHƯƠNG ÁN LỰA CHỌN
% ------------------------------------------------------------------------------
\noindent\textbf{PHẦN I. Câu trắc nghiệm nhiều phương án lựa chọn (3,0 điểm).}\\
\textit{Thí sinh trả lời từ câu 1 đến câu 12. Mỗi câu hỏi thí sinh chỉ chọn một phương án.}

\setcounter{ex}{0}

\begin{ex}%Câu 1
	Cho tập hợp $B = \{x \in \mathbb{N} \mid x < 4\}$. Viết tập hợp $B$ bằng cách liệt kê các phần tử, ta được
	\choice
	{$B = \{1; 2; 3\}$}
	{\True $B = \{0; 1; 2; 3\}$}
	{$B = \{0; 1; 2; 3; 4\}$}
	{$B = \{1; 2; 3; 4\}$}
	\loigiai{
		Các số tự nhiên nhỏ hơn $4$ gồm có $0; 1; 2; 3$, do đó $B = \{0; 1; 2; 3\}$.
	}
\end{ex}

\begin{ex}%Câu 2
	Số La Mã $\mathrm{XIV}$ có giá trị là
	\choice
	{$16$}
	{$9$}
	{$104$}
	{\True $14$}
	\loigiai{
		Trong hệ chữ số La Mã: $\mathrm{X} = 10$, $\mathrm{IV} = 4$. Do đó $\mathrm{XIV} = 10 + 4 = 14$.
	}
\end{ex}

\begin{ex}%Câu 3
	Số nhỏ nhất trong tập hợp $\mathbb{N}^*$ là
	\choice
	{\True $1$}
	{$0$}
	{$2$}
	{Không có số nhỏ nhất}
	\loigiai{
		Tập hợp $\mathbb{N}^* = \{1; 2; 3; \dots\}$, số nhỏ nhất trong tập hợp này là $1$.
	}
\end{ex}

\begin{ex}%Câu 4
	Kết quả của phép tính $5^7 : 5^3$ viết dưới dạng một luỹ thừa là
	\choice
	{$5^{10}$}
	{$5^{21}$}
	{\True $5^4$}
	{$1^4$}
	\loigiai{
		Áp dụng quy tắc chia hai luỹ thừa cùng cơ số: $5^7 : 5^3 = 5^{7 - 3} = 5^4$.
	}
\end{ex}

\begin{ex}%Câu 5
	Giá trị của biểu thức $20 - 2 \cdot 3^2$ là
	\choice
	{$162$}
	{\True $2$}
	{$14$}
	{$4$}
	\loigiai{
		Thực hiện luỹ thừa trước: $20 - 2 \cdot 3^2 = 20 - 2 \cdot 9 = 20 - 18 = 2$.
	}
\end{ex}

\begin{ex}%Câu 6
	Số nào sau đây là bội của $7$?
	\choice
	{$15$}
	{$27$}
	{$50$}
	{\True $35$}
	\loigiai{
		Vì $35 = 7 \cdot 5$ nên $35$ chia hết cho $7$, tức là $35$ là bội của $7$.
	}
\end{ex}

\begin{ex}%Câu 7
	Số nào sau đây chia hết cho $9$?
	\choice
	{\True $3\,726$}
	{$3\,725$}
	{$2\,019$}
	{$1\,234$}
	\loigiai{
		Số $3\,726$ có tổng các chữ số là $3 + 7 + 2 + 6 = 18$ chia hết cho $9$ nên $3\,726 \ \vdots \ 9$.
	}
\end{ex}

\begin{ex}%Câu 8
	Số nào sau đây là hợp số?
	\choice
	{$2$}
	{$17$}
	{\True $27$}
	{$31$}
	\loigiai{
		Số $27$ ngoài $1$ và $27$ còn có các ước là $3$ và $9$, do đó $27$ là hợp số.
	}
\end{ex}

\begin{ex}%Câu 9
	Số nào sau đây là bội chung của $4$ và $6$?
	\choice
	{$8$}
	{\True $24$}
	{$18$}
	{$30$}
	\loigiai{
		Ta thấy $24$ vừa chia hết cho $4$, vừa chia hết cho $6$ nên $24 \in \text{BC}(4, 6)$.
	}
\end{ex}

\begin{ex}%Câu 10
	Kết quả phân tích số $180$ ra thừa số nguyên tố là
	\choice
	{$2 \cdot 3^2 \cdot 10$}
	{$2^2 \cdot 45$}
	{$4 \cdot 9 \cdot 5$}
	{\True $2^2 \cdot 3^2 \cdot 5$}
	\loigiai{
		Ta có $180 = 4 \cdot 9 \cdot 5 = 2^2 \cdot 3^2 \cdot 5$.
	}
\end{ex}

\begin{ex}%Câu 11
	\noindent
	\begin{minipage}[t]{0.70\linewidth}
		Cho hình thoi $ABCD$ có hai đường chéo cắt nhau tại $O$ (hình bên). Khẳng định nào sau đây đúng?
		\choice
		{\True $AC \perp BD$}
		{$AC = BD$}
		{$AB = AC$}
		{$OA = OB$}
	\end{minipage}\hfill
	\begin{minipage}[t]{0.26\linewidth}
		\centering\vspace{0pt}
		\includegraphics[width=\linewidth]{images/de2_cau11.png}
	\end{minipage}
	\loigiai{
		Trong hình thoi, hai đường chéo vuông góc với nhau tại trung điểm của mỗi đường. Do đó khẳng định đúng là $AC \perp BD$.
	}
\end{ex}

\begin{ex}%Câu 12
	\noindent
	\begin{minipage}[t]{0.70\linewidth}
		Hình thang cân $ABCD$ ở hình bên có đáy nhỏ $AB = 4\text{ cm}$, đáy lớn $CD = 10\text{ cm}$, chiều cao $AH = 6\text{ cm}$. Diện tích hình thang cân $ABCD$ là
		\choice
		{$84\text{ cm}^2$}
		{$20\text{ cm}^2$}
		{\True $42\text{ cm}^2$}
		{$240\text{ cm}^2$}
	\end{minipage}\hfill
	\begin{minipage}[t]{0.26\linewidth}
		\centering\vspace{0pt}
		\includegraphics[width=\linewidth]{images/de2_cau12.png}
	\end{minipage}
	\loigiai{
		Diện tích hình thang cân $ABCD$ là:
		\[ S = \frac{(AB + CD) \cdot AH}{2} = \frac{(4 + 10) \cdot 6}{2} = 42\text{ (cm}^2\text{)}. \]
	}
\end{ex}

\vspace{8pt}

% ------------------------------------------------------------------------------
% PHẦN II: CÂU TRẮC NGHIỆM ĐÚNG SAI
% ------------------------------------------------------------------------------
\noindent\textbf{PHẦN II. Câu trắc nghiệm đúng sai (2,0 điểm).}\\
\textit{Thí sinh trả lời từ câu 13 đến câu 14. Trong mỗi ý a), b), c), d) ở mỗi câu, thí sinh chọn đúng hoặc sai.}

\begin{ex}%Câu 13
	Cho hai số $a = 54$ và $b = 90$.
	\choiceTF
	{$a$ chia hết cho $5$}
	{\True $3$ là ước chung của $a$ và $b$}
	{\True $\text{BCNN}(a, b) = 270$}
	{$a$ và $b$ có đúng $4$ ước chung}
	\loigiai{
		\begin{itemchoice}
			\item Sai: $54$ có chữ số tận cùng là $4$ nên không chia hết cho $5$.
			\item Đúng: Tổng các chữ số: $5 + 4 = 9$ và $9 + 0 = 9$ đều chia hết cho $3$, do đó $3$ là ước chung của $a$ và $b$.
			\item Đúng: Phân tích ra thừa số nguyên tố: $54 = 2 \cdot 3^3$ và $90 = 2 \cdot 3^2 \cdot 5$. Suy ra $\text{BCNN}(a, b) = 2 \cdot 3^3 \cdot 5 = 270$.
			\item Sai: Ta có $\text{ƯCLN}(a, b) = 2 \cdot 3^2 = 18$. Các ước chung của $a$ và $b$ là các ước của $18$: $\text{ƯC}(a, b) = \{1; 2; 3; 6; 9; 18\}$, gồm có $6$ phần tử, không phải $4$.
		\end{itemchoice}
	}
\end{ex}

\begin{ex}%Câu 14
	\noindent
	\begin{minipage}[t]{0.68\linewidth}
		Cho hình lục giác đều $ABCDEF$ có cạnh $4\text{ cm}$, tâm $O$ (hình bên).
		\choiceTF
		{\True Sáu góc của hình lục giác đều $ABCDEF$ bằng nhau}
		{$AD = AB$}
		{\True Chu vi hình lục giác đều $ABCDEF$ bằng $24\text{ cm}$}
		{Biết diện tích tam giác $OAB$ khoảng $6{,}9\text{ cm}^2$. Khi đó diện tích hình lục giác đều $ABCDEF$ khoảng $20{,}7\text{ cm}^2$}
	\end{minipage}\hfill
	\begin{minipage}[t]{0.28\linewidth}
		\centering\vspace{0pt}
		\includegraphics[width=\linewidth]{images/de2_cau14.png}
	\end{minipage}
	\loigiai{
		\begin{itemchoice}
			\item Đúng: Theo định nghĩa hình lục giác đều, sáu cạnh bằng nhau và sáu góc bằng nhau.
			\item Sai: Đường chéo chính $AD = OA + OD = 4 + 4 = 8\text{ cm}$, gấp đôi độ dài cạnh $AB = 4\text{ cm}$.
			\item Đúng: Chu vi hình lục giác đều bằng $6$ lần độ dài cạnh: $6 \cdot 4 = 24\text{ cm}$.
			\item Sai: Hình lục giác đều $ABCDEF$ được ghép bởi $6$ tam giác đều bằng nhau với tam giác $OAB$, do đó diện tích của lục giác đều khoảng $6 \cdot 6{,}9 = 41{,}4\text{ cm}^2$, không phải $20{,}7\text{ cm}^2$.
		\end{itemchoice}
	}
\end{ex}

\vspace{8pt}

% ------------------------------------------------------------------------------
% PHẦN III: CÂU TRẮC NGHIỆM TRẢ LỜI NGẮN
% ------------------------------------------------------------------------------
\noindent\textbf{PHẦN III. Câu trắc nghiệm trả lời ngắn (2,0 điểm).}\\
\textit{Thí sinh trả lời từ câu 15 đến câu 18.}

\begin{ex}%Câu 15
	Tính giá trị của biểu thức $3^4 - 4^3$.
	\shortans{17}
	\loigiai{
		$3^4 - 4^3 = 81 - 64 = 17$.
	}
\end{ex}

\begin{ex}%Câu 16
	Tìm chữ số $x$ để số $\overline{4x7}$ chia hết cho $9$.
	\shortans{7}
	\loigiai{
		Tổng các chữ số của số $\overline{4x7}$ là $4 + x + 7 = 11 + x$. Để số này chia hết cho $9$ thì $(11 + x) \ \vdots \ 9$.\\
		Vì $x$ là chữ số ($0 \le x \le 9$) nên $11 \le 11 + x \le 20$. Suy ra $11 + x = 18 \implies x = 7$.
	}
\end{ex}

\begin{ex}%Câu 17
	\noindent
	\begin{minipage}[t]{0.70\linewidth}
		Hình bình hành $ABCD$ ở hình bên có $AB = 8\text{ cm}$, $BC = 5\text{ cm}$. Tính chu vi hình bình hành $ABCD$ (đơn vị: cm).
		\shortans{26}
	\end{minipage}\hfill
	\begin{minipage}[t]{0.26\linewidth}
		\centering\vspace{0pt}
		\includegraphics[width=\linewidth]{images/de2_cau17.png}
	\end{minipage}
	\loigiai{
		Chu vi hình bình hành $ABCD$ là:
		\[ C = (AB + BC) \cdot 2 = (8 + 5) \cdot 2 = 26\text{ (cm)}. \]
	}
\end{ex}

\begin{ex}%Câu 18
	Bác Hà mua $4\text{ kg}$ táo, giá $35\,000$ đồng một ki-lô-gam và $3\text{ kg}$ cam, giá $28\,000$ đồng một ki-lô-gam. Bác đưa cho người bán $300\,000$ đồng. Hỏi người bán phải trả lại bác Hà bao nhiêu nghìn đồng?
	\shortans{76}
	\loigiai{
		Tổng số tiền mua táo và cam của bác Hà là:
		\[ 4 \cdot 35\,000 + 3 \cdot 28\,000 = 140\,000 + 84\,000 = 224\,000\text{ (đồng)}. \]
		Số tiền người bán phải trả lại cho bác Hà là:
		\[ 300\,000 - 224\,000 = 76\,000\text{ (đồng)}, \]
		tức là $76$ nghìn đồng.
	}
\end{ex}

\vspace{8pt}

% ------------------------------------------------------------------------------
% PHẦN IV: TỰ LUẬN
% ------------------------------------------------------------------------------
\noindent\textbf{PHẦN IV. Tự luận (3,0 điểm).}

\setcounter{bt}{18}

\begin{bt}%Câu 19
	(1,0 điểm).
	\begin{enumerate}[a)]
		\item Tính một cách hợp lí: $53 \cdot 125 - 53 \cdot 25 - 300$.
		\item Tìm số tự nhiên $x$, biết $5 \cdot 2^{x-1} + 12 = 92$.
	\end{enumerate}
	\begin{loigiai}
		\begin{enumerate}[a)]
			\item Ta có:
			\begin{align*}
				53 \cdot 125 - 53 \cdot 25 - 300 &= 53 \cdot (125 - 25) - 300\\
				&= 53 \cdot 100 - 300\\
				&= 5\,300 - 300 = 5\,000.
			\end{align*}
			\item Ta có:
			\begin{align*}
				5 \cdot 2^{x-1} + 12 &= 92\\
				5 \cdot 2^{x-1} &= 92 - 12\\
				5 \cdot 2^{x-1} &= 80\\
				2^{x-1} &= 80 : 5\\
				2^{x-1} &= 16\\
				2^{x-1} &= 2^4\\
				x - 1 &= 4\\
				x &= 5.
			\end{align*}
			Vậy $x = 5$.
		\end{enumerate}
	\end{loigiai}
\end{bt}

\begin{bt}%Câu 20
	(1,0 điểm).
	Số học sinh khối 6 của một trường trong khoảng từ $300$ đến $400$ em. Khi xếp thành hàng $12$, hàng $15$ hay hàng $18$ thì đều vừa đủ, không thừa em nào.
	\begin{enumerate}[a)]
		\item Tìm $\text{BCNN}(12, 15, 18)$.
		\item Tính số học sinh khối 6 của trường đó.
	\end{enumerate}
	\begin{loigiai}
		\begin{enumerate}[a)]
			\item Phân tích các số ra thừa số nguyên tố:
			\begin{align*}
				12 &= 2^2 \cdot 3;\\
				15 &= 3 \cdot 5;\\
				18 &= 2 \cdot 3^2.
			\end{align*}
			Do đó:
			\[ \text{BCNN}(12, 15, 18) = 2^2 \cdot 3^2 \cdot 5 = 4 \cdot 9 \cdot 5 = 180. \]
			\item Gọi số học sinh khối 6 của trường đó là $x$ ($x \in \mathbb{N}^*$, $300 \le x \le 400$).\\
			Theo đề bài, khi xếp thành hàng $12$, hàng $15$, hàng $18$ đều vừa đủ nên $x$ chia hết cho cả $12, 15, 18$.\\
			Suy ra $x \in \text{BC}(12, 15, 18)$.\\
			Vì $\text{BCNN}(12, 15, 18) = 180$ nên:
			\[ \text{BC}(12, 15, 18) = \text{B}(180) = \{0; 180; 360; 540; \dots\}. \]
			Mà $300 \le x \le 400$ nên $x = 360$.\\
			Vậy khối 6 của trường đó có $360$ học sinh.
		\end{enumerate}
	\end{loigiai}
\end{bt}

\begin{bt}%Câu 21
	(1,0 điểm).
	
	\noindent
	\begin{minipage}[t]{0.68\linewidth}
		Nền một căn phòng có dạng hình chữ nhật $ABCD$ với chiều dài $AB = 8\text{ m}$, chiều rộng $BC = 6\text{ m}$ (hình bên).
		\begin{enumerate}[a)]
			\item Tính chu vi và diện tích nền căn phòng.
			\item Người ta lát kín nền căn phòng bằng các viên gạch hình vuông cạnh $40\text{ cm}$, giá $25\,000$ đồng một viên. Tính số tiền mua gạch (coi như mạch vữa không đáng kể và không có viên gạch nào bị hỏng).
		\end{enumerate}
	\end{minipage}\hfill
	\begin{minipage}[t]{0.28\linewidth}
		\centering\vspace{0pt}
		\includegraphics[width=\linewidth]{images/de2_cau21.png}
	\end{minipage}
	\begin{loigiai}
		\begin{enumerate}[a)]
			\item Chu vi nền căn phòng là:
			\[ (8 + 6) \cdot 2 = 28\text{ (m)}. \]
			Diện tích nền căn phòng là:
			\[ 8 \cdot 6 = 48\text{ (m}^2\text{)}. \]
			\item Đổi: $8\text{ m} = 800\text{ cm}$; $6\text{ m} = 600\text{ cm}$.\\
			Số viên gạch được lát theo chiều dài căn phòng là:
			\[ 800 : 40 = 20\text{ (viên)}. \]
			Số viên gạch được lát theo chiều rộng căn phòng là:
			\[ 600 : 40 = 15\text{ (hàng)}. \]
			Tổng số viên gạch cần dùng để lát kín nền phòng là:
			\[ 20 \cdot 15 = 300\text{ (viên)}. \]
			\textit{(Cách khác: Diện tích một viên gạch là $40 \cdot 40 = 1\,600\text{ cm}^2 = 0{,}16\text{ m}^2$. Số viên gạch cần dùng là $48 : 0{,}16 = 300$ viên).}\\
			Số tiền mua gạch để lát nền căn phòng là:
			\[ 300 \cdot 25\,000 = 7\,500\,000\text{ (đồng)}. \]
		\end{enumerate}
	\end{loigiai}
\end{bt}

\begin{center}
	\rule{4cm}{0.5pt} \textbf{HẾT} \rule{4cm}{0.5pt}\\
	\textit{(Giáo viên coi kiểm tra không giải thích gì thêm.)}
\end{center}

\newpage
% ==============================================================================
% BẢNG ĐÁP ÁN VÀ HƯỚNG DẪN CHẤM - ĐỀ 2
% ==============================================================================
\begin{center}
	{\Large\textbf{ĐÁP ÁN VÀ HƯỚNG DẪN CHẤM -- ĐỀ 2}}\\[4pt]
	\textbf{(Toán 6, giữa học kì I, 2026 -- 2027)}
\end{center}

\vspace{6pt}

\noindent\textbf{PHẦN I. (3,0 điểm) Mỗi câu trả lời đúng được 0,25 điểm.}

\begin{center}
\renewcommand{\arraystretch}{1.3}
\begin{tabular}{|c|c|c|c|c|c|c|c|c|c|c|c|c|}
	\hline
	\textbf{Câu} & \textbf{1} & \textbf{2} & \textbf{3} & \textbf{4} & \textbf{5} & \textbf{6} & \textbf{7} & \textbf{8} & \textbf{9} & \textbf{10} & \textbf{11} & \textbf{12} \\
	\hline
	\textbf{Đáp án} & B & D & A & C & B & D & A & C & B & D & A & C \\
	\hline
\end{tabular}
\end{center}

\vspace{4pt}

\noindent\textbf{PHẦN II. (2,0 điểm)}\\
\textit{Điểm tối đa của mỗi câu là 1 điểm: chọn đúng 1 ý được 0,1 điểm; đúng 2 ý được 0,25 điểm; đúng 3 ý được 0,5 điểm; đúng cả 4 ý được 1 điểm.}

\begin{center}
\renewcommand{\arraystretch}{1.3}
\begin{tabular}{|c|c|c|p{11.5cm}|}
	\hline
	\textbf{Câu} & \textbf{Ý} & \textbf{Đáp án} & \centering\textbf{Giải thích ngắn} \tabularnewline
	\hline
	& a) & \textbf{Sai} & $54$ có chữ số tận cùng là $4$ nên không chia hết cho $5$. \tabularnewline
	\cline{2-4}
	& b) & \textbf{Đúng} & $5 + 4 = 9$ và $9 + 0 = 9$ đều chia hết cho $3$ nên $3$ là ước chung của $a$ và $b$. \tabularnewline
	\cline{2-4}
	\textbf{13} & c) & \textbf{Đúng} & $54 = 2 \cdot 3^3$; $90 = 2 \cdot 3^2 \cdot 5$ nên $\text{BCNN}(a, b) = 2 \cdot 3^3 \cdot 5 = 270$. \tabularnewline
	\cline{2-4}
	& d) & \textbf{Sai} & $\text{ƯCLN}(a, b) = 2 \cdot 3^2 = 18$; $\text{ƯC}(a, b) = \text{Ư}(18) = \{1; 2; 3; 6; 9; 18\}$ có $6$ phần tử, không phải $4$. \tabularnewline
	\hline
	& a) & \textbf{Đúng} & Hình lục giác đều có sáu cạnh bằng nhau và sáu góc bằng nhau. \tabularnewline
	\cline{2-4}
	& b) & \textbf{Sai} & Đường chéo chính $AD = OA + OD = 4 + 4 = 8\text{ (cm)}$, gấp đôi cạnh $AB$. \tabularnewline
	\cline{2-4}
	\textbf{14} & c) & \textbf{Đúng} & Chu vi: $6 \cdot 4 = 24\text{ (cm)}$. \tabularnewline
	\cline{2-4}
	& d) & \textbf{Sai} & Hình lục giác đều ghép bởi $6$ tam giác đều bằng nhau như $OAB$ nên diện tích khoảng $6 \cdot 6{,}9 = 41{,}4\text{ (cm}^2\text{)}$, không phải $20{,}7\text{ cm}^2$. \tabularnewline
	\hline
\end{tabular}
\end{center}

\vspace{4pt}

\noindent\textbf{PHẦN III. (2,0 điểm) Mỗi câu trả lời đúng được 0,5 điểm.}

\begin{center}
\renewcommand{\arraystretch}{1.3}
\begin{tabular}{|c|c|p{12.5cm}|}
	\hline
	\textbf{Câu} & \textbf{Đáp án} & \centering\textbf{Lời giải tóm tắt} \tabularnewline
	\hline
	\textbf{15} & \textbf{17} & $3^4 - 4^3 = 81 - 64 = 17$. \tabularnewline
	\hline
	\textbf{16} & \textbf{7} & $4 + x + 7 = 11 + x$ chia hết cho $9$, mà $x$ là chữ số ($0 \le x \le 9$) nên $11 + x = 18 \implies x = 7$. \tabularnewline
	\hline
	\textbf{17} & \textbf{26} & Chu vi hình bình hành: $(8 + 5) \cdot 2 = 26\text{ (cm)}$. \tabularnewline
	\hline
	\textbf{18} & \textbf{76} & Số tiền phải trả: $4 \cdot 35\,000 + 3 \cdot 28\,000 = 224\,000\text{ (đồng)}$. Tiền trả lại: $300\,000 - 224\,000 = 76\,000\text{ (đồng)}$, tức $76$ nghìn đồng. \tabularnewline
	\hline
\end{tabular}
\end{center}

\vspace{4pt}

\noindent\textbf{PHẦN IV. Tự luận (3,0 điểm)}\\
\textit{Hình vẽ Câu 21 cho sẵn trong đề, không tính điểm vẽ hình. Cách giải khác đúng vẫn cho điểm tối đa.}

\begin{center}
\renewcommand{\arraystretch}{1.3}
\begin{tabular}{|c|l|c|}
	\hline
	\textbf{Câu} & \centering\textbf{Nội dung} & \textbf{Điểm} \tabularnewline
	\hline
	\textbf{19 a)} & $53 \cdot 125 - 53 \cdot 25 - 300 = 53 \cdot (125 - 25) - 300$ & 0,25 \tabularnewline
	& $= 53 \cdot 100 - 300 = 5\,300 - 300 = 5\,000$ & 0,25 \tabularnewline
	\hline
	\textbf{19 b)} & $5 \cdot 2^{x-1} = 92 - 12 = 80 \implies 2^{x-1} = 16$ & 0,25 \tabularnewline
	& $2^{x-1} = 2^4 \implies x - 1 = 4 \implies x = 5$ & 0,25 \tabularnewline
	\hline
	& $12 = 2^2 \cdot 3$; $15 = 3 \cdot 5$; $18 = 2 \cdot 3^2$. & 0,25 \tabularnewline
	\cline{2-3}
	\textbf{20 a)} & $\text{BCNN}(12, 15, 18) = 2^2 \cdot 3^2 \cdot 5 = 180$. & 0,25 \tabularnewline
	\hline
	& Gọi số học sinh là $x$. Vì xếp hàng $12, 15, 18$ đều vừa đủ nên $x \in \text{BC}(12, 15, 18)$ & \tabularnewline
	\textbf{20 b)} & $= \text{B}(180) = \{0; 180; 360; 540; \dots\}$. & 0,25 \tabularnewline
	\cline{2-3}
	& Mà $300 \le x \le 400$ nên $x = 360$. Khối 6 có $360$ học sinh. & 0,25 \tabularnewline
	\hline
	& Chu vi nền căn phòng: $(8 + 6) \cdot 2 = 28\text{ (m)}$. & 0,25 \tabularnewline
	\cline{2-3}
	\textbf{21 a)} & Diện tích nền căn phòng: $8 \cdot 6 = 48\text{ (m}^2\text{)}$. & 0,25 \tabularnewline
	\hline
	& Đổi $8\text{ m} = 800\text{ cm}$, $6\text{ m} = 600\text{ cm}$. & \tabularnewline
	\textbf{21 b)} & Mỗi hàng theo chiều dài có $800 : 40 = 20$ viên, theo chiều rộng có $600 : 40 = 15$ hàng; & \tabularnewline
	& số viên gạch: $20 \cdot 15 = 300$ (viên). & 0,25 \tabularnewline
	\cline{2-3}
	& Số tiền mua gạch: $300 \cdot 25\,000 = 7\,500\,000\text{ (đồng)}$. & 0,25 \tabularnewline
	\hline
\end{tabular}
\end{center}
